%% ============================================================
\section*{Metric Convention Summary}
%% ============================================================

All formulas in this document are derived in \textbf{both} the West coast and East coast
metric conventions. The table below summarises the key differences:

\medskip
\begin{center}
\renewcommand{\arraystretch}{1.4}
\begin{tabular}{@{}lcc@{}}
  \toprule
  Property & West coast $(+{-}{-}{-})$ & East coast $(-{+}{+}{+})$ \\
  \midrule
  $g^{\mu\nu}$ & $\mathrm{diag}(+1,-1,-1,-1)$ & $\mathrm{diag}(-1,+1,+1,+1)$ \\
  $g^{00}$ & $+1$ & $-1$ \\
  $g^{11}$ & $-1$ & $+1$ \\
  Signature & $+2$ & $-2$ \\
  Used by & Srednicki, Weinberg & Peskin--Schroeder, most particle physicists \\
  \bottomrule
\end{tabular}
\end{center}
\medskip

\noindent\textbf{Key consequence}: any formula involving $g^{\mu\nu}$ (boosted generators,
$\sigma^\mu$, trace identities) differs by signs between conventions.
Formulas involving only spatial indices ($g^{ij}$) pick up a sign flip
($-1$ in WC vs.\ $+1$ in EC), while the time component $g^{00}$ flips
($+1$ in WC vs.\ $-1$ in EC).

\tableofcontents
\newpage

%% ============================================================
\section{The Lorentz Algebra and its
  \texorpdfstring{$\mathfrak{su}(2)\oplus\mathfrak{su}(2)$}{su(2)+su(2)}
  Structure}
%% ============================================================

\subsection{The six generators of \texorpdfstring{$\mathrm{SO}(1,3)$}{SO(1,3)}}

The Lorentz group $\mathrm{SO}(1,3)$ is a six-dimensional Lie group.
Its Lie algebra $\mathfrak{so}(1,3)$ is spanned by six generators,
conventionally split into:
\begin{itemize}
  \item \textbf{Rotations:} $J_i = \tfrac{1}{2}\varepsilon_{ijk}M^{jk}$,\quad $i=1,2,3$
  \item \textbf{Boosts:} $K_i = M^{i0}$,\quad $i=1,2,3$
\end{itemize}
where $M^{\mu\nu} = -M^{\nu\mu}$ are the antisymmetric generators of the
defining (4-dimensional vector) representation.

The commutation relations of the Lorentz algebra are:
\begin{align}
  [J_i,\, J_j] &= i\,\varepsilon_{ijk}\,J_k
  \label{eq:JJ}\\
  [J_i,\, K_j] &= i\,\varepsilon_{ijk}\,K_k
  \label{eq:JK}\\
  [K_i,\, K_j] &= -i\,\varepsilon_{ijk}\,J_k
  \label{eq:KK}
\end{align}
Equation~\eqref{eq:JJ} says the rotations by themselves close into an
$\mathfrak{su}(2)$ algebra.
Equation~\eqref{eq:JK} says the boosts $K_i$ transform as a 3-vector under
rotations --- exactly what one expects from a physical boost parameter.
Equation~\eqref{eq:KK} is the key relation: two boosts do \emph{not} close
into another boost; instead they produce a rotation (Thomas precession).
The crucial minus sign in~\eqref{eq:KK} is a direct consequence of the
Lorentzian signature: in $\mathrm{SO}(4)$ (Euclidean) one would get $+i\varepsilon_{ijk}J_k$.

\begin{remark}
  Equations \eqref{eq:JJ}--\eqref{eq:KK} are convention-independent:
  they hold in both the West coast $(+{-}{-}{-})$ and East coast
  $(-{+}{+}{+})$ signatures when $J_i$ and $K_i$ are defined as above.
  (One can verify this from the master commutator
  $[M^{\mu\nu},M^{\rho\sigma}] = i(g^{\nu\rho}M^{\mu\sigma}
  - g^{\mu\rho}M^{\nu\sigma} - g^{\nu\sigma}M^{\mu\rho}
  + g^{\mu\sigma}M^{\nu\rho})$
  by substituting the explicit definitions of $J_i$ and $K_i$.)
\end{remark}

%%------------------------------------------------------------
\subsection{The \texorpdfstring{$N_i,\,N_i^\dagger$}{N, N†} combinations}

Define new combinations (Srednicki eq.~33.16):
\begin{align}
  N_i &\equiv \tfrac{1}{2}(J_i - iK_i),
  \label{eq:Ndef}\\
  N_i^\dagger &\equiv \tfrac{1}{2}(J_i + iK_i).
  \label{eq:Ndagdef}
\end{align}
We claim these satisfy
\begin{align}
  [N_i, N_j] &= i\varepsilon_{ijk}N_k,
  \label{eq:NN}\\
  [N_i^\dagger, N_j^\dagger] &= i\varepsilon_{ijk}N_k^\dagger,
  \label{eq:NdagNdag}\\
  [N_i, N_j^\dagger] &= 0.
  \label{eq:NNdag}
\end{align}
We now derive each relation explicitly, step by step.

%%------------------------------------------------------------
\subsection{Derivation of \texorpdfstring{$[N_i,N_j]=i\varepsilon_{ijk}N_k$}{[N,N]=iepsN}}

\paragraph{Step 1: Expand the commutator.}
\begin{align}
  [N_i, N_j]
  &= \Bigl[\tfrac{1}{2}(J_i - iK_i),\; \tfrac{1}{2}(J_j - iK_j)\Bigr]
  \notag\\
  &= \tfrac{1}{4}\Bigl[J_i - iK_i,\; J_j - iK_j\Bigr].
\end{align}
Using bilinearity of the commutator:
\begin{align}
  [J_i - iK_i,\; J_j - iK_j]
  &= [J_i, J_j]
     - i[J_i, K_j]
     - i[K_i, J_j]
     + i^2[K_i, K_j]
  \notag\\
  &= [J_i, J_j]
     - i[J_i, K_j]
     - i[K_i, J_j]
     - [K_i, K_j].
  \label{eq:NN_expand}
\end{align}

\paragraph{Step 2: Substitute the three Lorentz commutators.}
\begin{align*}
  [J_i, J_j] &= i\varepsilon_{ijk}J_k,\\
  [J_i, K_j] &= i\varepsilon_{ijk}K_k,\\
  [K_i, J_j] &= -[J_j, K_i] = -i\varepsilon_{jik}K_k = +i\varepsilon_{ijk}K_k,\\
  [K_i, K_j] &= -i\varepsilon_{ijk}J_k.
\end{align*}
(In the third line we used $\varepsilon_{jik} = -\varepsilon_{ijk}$.)

\paragraph{Step 3: Substitute into \eqref{eq:NN_expand}.}
\begin{align}
  [J_i - iK_i,\; J_j - iK_j]
  &= i\varepsilon_{ijk}J_k
     - i(i\varepsilon_{ijk}K_k)
     - i(i\varepsilon_{ijk}K_k)
     - (-i\varepsilon_{ijk}J_k)
  \notag\\
  &= i\varepsilon_{ijk}J_k
     + \varepsilon_{ijk}K_k
     + \varepsilon_{ijk}K_k
     + i\varepsilon_{ijk}J_k
  \notag\\
  &= 2i\varepsilon_{ijk}J_k + 2\varepsilon_{ijk}K_k
  \notag\\
  &= 2\varepsilon_{ijk}(iJ_k + K_k).
  \label{eq:step3}
\end{align}

\paragraph{Step 4: Identify the result.}
Note that $iJ_k + K_k = i(J_k - iK_k)$. Therefore:
\begin{align}
  [J_i - iK_i,\; J_j - iK_j]
  &= 2\varepsilon_{ijk}\cdot i(J_k - iK_k)
   = 4i\varepsilon_{ijk}\cdot\tfrac{1}{2}(J_k - iK_k)
   = 4i\varepsilon_{ijk}N_k.
\end{align}
Multiplying by $\tfrac{1}{4}$:
\begin{equation}
  \boxed{[N_i, N_j] = \tfrac{1}{4}\cdot 4i\varepsilon_{ijk}N_k = i\varepsilon_{ijk}N_k.}
\end{equation}

\paragraph{Key insight on signs.}
The crucial role of the sign in $[K_i,K_j]=-i\varepsilon_{ijk}J_k$ is
visible in Step~3: the $i^2 = -1$ from the $(-i)(-i)$ prefactor
\emph{combined} with the minus sign in $[K_i,K_j]$
produces a \emph{double negative}, giving $+i\varepsilon_{ijk}J_k$
contribution that adds constructively with the $[J_i,J_j]$ term.
Had we instead had $[K_i,K_j]=+i\varepsilon_{ijk}J_k$
(as in $\mathfrak{so}(4)$), the $J$ terms would cancel and we would not
obtain a closed $\mathfrak{su}(2)$ algebra for $N_i$.

%%------------------------------------------------------------
\subsection{Derivation of
  \texorpdfstring{$[N_i^\dagger,N_j^\dagger]=i\varepsilon_{ijk}N_k^\dagger$}%
    {[N†,N†]=iepsN†}}

By an identical calculation with $iK \to -iK$:
\begin{align}
  [J_i + iK_i,\; J_j + iK_j]
  &= [J_i, J_j]
     + i[J_i, K_j]
     + i[K_i, J_j]
     + i^2[K_i, K_j]
  \notag\\
  &= i\varepsilon_{ijk}J_k
     + i(i\varepsilon_{ijk}K_k)
     + i(i\varepsilon_{ijk}K_k)
     + (-1)(-i\varepsilon_{ijk}J_k)
  \notag\\
  &= i\varepsilon_{ijk}J_k
     - \varepsilon_{ijk}K_k
     - \varepsilon_{ijk}K_k
     + i\varepsilon_{ijk}J_k
  \notag\\
  &= 2i\varepsilon_{ijk}J_k - 2\varepsilon_{ijk}K_k
  \notag\\
  &= 2\varepsilon_{ijk}(iJ_k - K_k)
   = 4i\varepsilon_{ijk}\cdot\tfrac{1}{2}(J_k + iK_k)
   = 4i\varepsilon_{ijk}N_k^\dagger.
\end{align}
Hence:
\begin{equation}
  \boxed{[N_i^\dagger, N_j^\dagger] = i\varepsilon_{ijk}N_k^\dagger.}
\end{equation}

%%------------------------------------------------------------
\subsection{Derivation of
  \texorpdfstring{$[N_i,N_j^\dagger]=0$}{[N,N†]=0}}

\begin{align}
  [N_i, N_j^\dagger]
  &= \Bigl[\tfrac{1}{2}(J_i - iK_i),\; \tfrac{1}{2}(J_j + iK_j)\Bigr]
  \notag\\
  &= \tfrac{1}{4}\Bigl[J_i - iK_i,\; J_j + iK_j\Bigr]
  \notag\\
  &= \tfrac{1}{4}\Bigl(
       [J_i, J_j] + i[J_i, K_j] - i[K_i, J_j] - i^2[K_i, K_j]
     \Bigr)
  \notag\\
  &= \tfrac{1}{4}\Bigl(
       i\varepsilon_{ijk}J_k
       + i(i\varepsilon_{ijk}K_k)
       - i(i\varepsilon_{ijk}K_k)
       + [K_i, K_j]
     \Bigr)
  \notag\\
  &= \tfrac{1}{4}\Bigl(
       i\varepsilon_{ijk}J_k
       - \varepsilon_{ijk}K_k
       + \varepsilon_{ijk}K_k
       + (-i\varepsilon_{ijk}J_k)
     \Bigr)
  \notag\\
  &= \tfrac{1}{4}\Bigl(i\varepsilon_{ijk}J_k - i\varepsilon_{ijk}J_k\Bigr) = 0.
\end{align}
\begin{equation}
  \boxed{[N_i, N_j^\dagger] = 0.}
\end{equation}

%%------------------------------------------------------------
\subsection{The isomorphism statement --- over \texorpdfstring{$\mathbb{C}$}{C}}

The computation above shows that the six operators
$\{N_i, N_i^\dagger\}$ split into two commuting copies of the
$\mathfrak{su}(2)$ algebra.  However, one must be careful about the field:

\begin{remark}[Over $\mathbb{R}$: no isomorphism]
  Over the reals, $\mathfrak{so}(1,3) \not\cong \mathfrak{su}(2)\oplus\mathfrak{su}(2)$.
  The reason is that the generators $K_i$ of boosts are \emph{not Hermitian}
  (and cannot be made so while keeping $J_i$ Hermitian, because the Lorentz
  group is non-compact).  Consequently $N_i^\dagger$ is the Hermitian
  conjugate of $N_i$, not an independent generator: there are only 6
  real independent generators, not 6+6.
  The real form $\mathfrak{su}(2)\oplus\mathfrak{su}(2)$ has only
  \emph{compact} unitary representations, while $\mathrm{SO}(1,3)$ has
  no finite-dimensional unitary representations (other than the trivial one).
\end{remark}

\begin{remark}[Over $\mathbb{C}$: the isomorphism holds]
  Consider the \emph{complexification}
  $\mathfrak{so}(1,3)_\mathbb{C} \equiv \mathfrak{so}(1,3)\otimes_\mathbb{R}\mathbb{C}$.
  This is the complex Lie algebra obtained by allowing complex linear combinations
  of the real generators $\{J_i, K_i\}$.  In this larger algebra, $N_i$ and
  $N_i^\dagger$ \emph{are} independent (not related by a reality condition),
  and our calculation proves:
  \begin{equation}
    \mathfrak{so}(1,3)_\mathbb{C}
    \;\cong\;
    \mathfrak{sl}(2,\mathbb{C}) \oplus \mathfrak{sl}(2,\mathbb{C})
    \;\cong\;
    \mathfrak{su}(2)_\mathbb{C} \oplus \mathfrak{su}(2)_\mathbb{C}.
    \label{eq:iso}
  \end{equation}
  The two copies are called $\mathfrak{su}(2)_L$ (generated by $N_i$) and
  $\mathfrak{su}(2)_R$ (generated by $N_i^\dagger$).
  (The subscripts $L/R$ correspond to left-handed and right-handed chirality
  of the associated spinors.)
\end{remark}

The physical content of~\eqref{eq:iso} is that the representation theory
of the Lorentz group (for finite-dimensional, not necessarily unitary, reps)
reduces to \emph{two independent} $\mathfrak{su}(2)$ problems.

%%------------------------------------------------------------
\subsection{Irreducible representations: the \texorpdfstring{$(n,n')$}{(n,n')} labels}

Because $\mathfrak{su}(2)_L \oplus \mathfrak{su}(2)_R$ is semisimple with
each factor being $\mathfrak{su}(2)$, its irreducible representations
are tensor products:
\begin{equation}
  V_{(n,n')} = V_n^L \otimes V_{n'}^R,
\end{equation}
where $V_n^L$ is the spin-$n$ irrep of $\mathfrak{su}(2)_L$ and
$V_{n'}^R$ is the spin-$n'$ irrep of $\mathfrak{su}(2)_R$,
with $n,n' \in \{0, \tfrac{1}{2}, 1, \tfrac{3}{2}, \ldots\}$.

The dimension of the $(n,n')$ representation is:
\begin{equation}
  \dim(n,n') = (2n+1)(2n'+1).
\end{equation}

\medskip
\noindent\textbf{Table of physical representations.}  In Srednicki's
notation, representations are often labelled by their dimensions
$(2n+1, 2n'+1)$:

\begin{center}
\renewcommand{\arraystretch}{1.45}
\begin{tabular}{@{}ccccl@{}}
  \toprule
  $(n,n')$ & Srednicki dim & dim & Physical field & Notes \\
  \midrule
  $(0,0)$ & $(1,1)$ & 1 & Scalar $\phi$ & Trivial rep \\
  $(\tfrac{1}{2},0)$ & $(2,1)$ & 2 & Left Weyl $\psi_a$ & $N_i^\dagger$ trivial \\
  $(0,\tfrac{1}{2})$ & $(1,2)$ & 2 & Right Weyl $\psidag_{\dot a}$ & $N_i$ trivial \\
  $(\tfrac{1}{2},\tfrac{1}{2})$ & $(2,2)$ & 4 & 4-vector $A^\mu$ & \\
  $(1,0)$ & $(3,1)$ & 3 & Self-dual 2-form $F^+_{\mu\nu}$ & Left field strength \\
  $(0,1)$ & $(1,3)$ & 3 & Anti-self-dual 2-form $F^-_{\mu\nu}$ & Right field strength \\
  $(1,1)$ & $(3,3)$ & 9 & Traceless symmetric tensor $h^{\mu\nu}_\perp$ & Graviton \\
  $(\tfrac{1}{2},1)$ & $(2,3)$ & 6 & Rarita-Schwinger $\psi^\mu_a$ & Gravitino \\
  \bottomrule
\end{tabular}
\end{center}

\begin{physinterp}
  The $(n,n')$ labels have direct physical meaning:
  $n$ is the ``left-handed spin'' (the eigenvalue of $|\vec{N}|$)
  and $n'$ is the ``right-handed spin'' (the eigenvalue of $|\vec{N}^\dagger|$).
  A field transforming in the $(\tfrac{1}{2},0)$ rep is a 2-component left-handed
  Weyl spinor because $\mathfrak{su}(2)_L$ acts nontrivially (as the spin-$\tfrac{1}{2}$
  representation) while $\mathfrak{su}(2)_R$ acts trivially.
  Under parity, $J_i \to J_i$ but $K_i \to -K_i$, which swaps $N_i \leftrightarrow N_i^\dagger$,
  hence $(n,n') \to (n',n)$.
  A left-handed Weyl spinor $(\tfrac{1}{2},0)$ maps to a right-handed one $(0,\tfrac{1}{2})$
  under parity --- exactly what we expect physically.
\end{physinterp}

%% ============================================================
\section{Tensor Product Decompositions of Lorentz Representations}
%% ============================================================

\subsection{Warmup: \texorpdfstring{$\mathrm{SU}(2)$}{SU(2)} Clebsch-Gordan decomposition}

For $\mathrm{SU}(2)$, representations are labelled by spin $j \in \{0,\tfrac{1}{2},1,\ldots\}$,
and the tensor product decomposes via the Clebsch-Gordan series:
\begin{equation}
  j_1 \otimes j_2 = |j_1-j_2| \oplus (|j_1-j_2|+1) \oplus \cdots \oplus (j_1+j_2).
  \label{eq:CG}
\end{equation}
The most important case is the fundamental representation $j=\tfrac{1}{2}$
(the 2-dimensional ``doublet'' or ``spinor'' representation):
\begin{equation}
  \tfrac{1}{2} \otimes \tfrac{1}{2}
  = 0 \oplus 1,
  \qquad\text{i.e., in dimensions: }\quad
  \mathbf{2} \otimes \mathbf{2} = \mathbf{1} \oplus \mathbf{3}.
\end{equation}

\paragraph{Why this decomposition?}
The tensor product space $\mathbf{2}\otimes\mathbf{2}$ is 4-dimensional.
Label the basis of each $\mathbf{2}$ as $\{|{\uparrow}\rangle, |{\downarrow}\rangle\}$,
with indices $a,b \in \{1,2\}$.  A general element of $\mathbf{2}\otimes\mathbf{2}$
is a tensor $T^{ab}$ with two indices.  We decompose into irreducible parts using
the standard Young tableau approach --- projecting onto definite symmetry under
$a \leftrightarrow b$:

\begin{itemize}
  \item \textbf{Antisymmetric subspace} (Young tableau: single column of 2 boxes):
  \[
    T^{[ab]} = \tfrac{1}{2}(T^{ab} - T^{ba}).
  \]
  This is a 1-dimensional space (for $2\times 2$ antisymmetric matrices).
  The unique (up to scale) invariant tensor is $\varepsilon^{ab}$, and the singlet state is:
  \[
    |0,0\rangle = \frac{|{\uparrow\downarrow}\rangle - |{\downarrow\uparrow}\rangle}{\sqrt{2}}
    = \frac{1}{\sqrt{2}}\varepsilon^{ab}|a\rangle\otimes|b\rangle.
  \]
  This is the spin-0 singlet, the $j=0$ rep. Its dimension is 1.

  \item \textbf{Symmetric subspace} (Young tableau: single row of 2 boxes):
  \[
    T^{(ab)} = \tfrac{1}{2}(T^{ab} + T^{ba}).
  \]
  This is a 3-dimensional space (for $2\times 2$ symmetric matrices).
  Basis: $|{\uparrow\uparrow}\rangle$,
  $(|{\uparrow\downarrow}\rangle+|{\downarrow\uparrow}\rangle)/\sqrt{2}$,
  $|{\downarrow\downarrow}\rangle$ --- the three states $|1,m\rangle$ for $m=1,0,-1$.
  This is the spin-1 triplet, the $j=1$ rep.  Its dimension is 3.
\end{itemize}

We write this compactly as:
\begin{equation}
  \mathbf{2}\otimes\mathbf{2} = \mathbf{1}_A \oplus \mathbf{3}_S,
  \label{eq:su2_22}
\end{equation}
where the subscripts $A$ (antisymmetric) and $S$ (symmetric) refer to
the symmetry of the tensor product under index exchange.

\paragraph{The invariant tensor $\varepsilon^{ab}$.}
The fact that a singlet ($j=0$, trivial rep) appears in $\mathbf{2}\otimes\mathbf{2}$
is equivalent to the existence of an invariant antisymmetric 2-index tensor
$\varepsilon^{ab}$ (with $\varepsilon^{12}=-\varepsilon^{21}=1$).
Under $\mathrm{SU}(2)$ transformations $U$: $T^{ab} \to U^a{}_c U^b{}_d T^{cd}$,
and $\varepsilon^{ab}$ is invariant because $\det U = 1$ for $U\in\mathrm{SU}(2)$.
This $\varepsilon^{ab}$ is the metric of the spinor space --- it raises and lowers
spinor indices just as $g^{\mu\nu}$ does for spacetime vectors.

%%------------------------------------------------------------
\subsection{Lorentz tensor products: \texorpdfstring{$(2,1)\otimes(2,1)$}{(2,1)x(2,1)}}

Now we work with the Lorentz group.  Using the isomorphism
$\mathrm{SO}(1,3)_\mathbb{C} \cong \mathrm{SU}(2)_L \times \mathrm{SU}(2)_R$,
a representation $(d_L, d_R)$ (in Srednicki's dimension notation) means:
$d_L = 2n_L+1$-dimensional under $\mathrm{SU}(2)_L$, and
$d_R = 2n_R+1$-dimensional under $\mathrm{SU}(2)_R$.

The tensor product decomposes factor by factor:
\begin{equation}
  (d_L, d_R) \otimes (d_L', d_R')
  = (d_L \otimes_L d_L') \otimes (d_R \otimes_R d_R'),
\end{equation}
where each factor is decomposed via the $\mathrm{SU}(2)$ Clebsch-Gordan rule.

\medskip
\noindent\textbf{Case: $(2,1)\otimes(2,1)$, i.e., $(\tfrac{1}{2},0)\otimes(\tfrac{1}{2},0)$.}

This is the tensor product of two left-handed Weyl spinors $\psi_a \otimes \chi_b$,
which lives in a space with indices $(a,b)$.

\begin{align*}
  \text{SU(2)}_L\text{ factor:} &\quad
  \tfrac{1}{2}\otimes\tfrac{1}{2} = 0_A \oplus 1_S,
  \quad\text{i.e.,}\quad \mathbf{2}\otimes\mathbf{2} = \mathbf{1}_A \oplus \mathbf{3}_S.\\[3pt]
  \text{SU(2)}_R\text{ factor:} &\quad
  0\otimes 0 = 0\quad\text{(trivial: stays 1-dimensional).}
\end{align*}

Combining:
\begin{equation}
  \boxed{(2,1)\otimes(2,1) = (1,1)_A \oplus (3,1)_S.}
  \label{eq:21x21}
\end{equation}
Dimension check: $2\times 2 = 1\times 1 + 3\times 1 = 1 + 3 = 4$. $\checkmark$

\paragraph{What is the $(1,1)_A$ piece?}
The antisymmetric part of $\psi_a\chi_b$ is:
\[
  \psi_a\chi_b - \psi_b\chi_a
  = -\varepsilon_{ab}\,(\varepsilon^{cd}\psi_c\chi_d)
  = -\varepsilon_{ab}\,\psi^c\chi_c.
\]
The scalar $\psi^c\chi_c = \psi^a\chi_a$ is a Lorentz scalar (the antisymmetric
singlet), and $\varepsilon_{ab}$ is its ``carrier'' --- the invariant tensor
of $\mathrm{SL}(2,\mathbb{C})$.  This is exactly the epsilon symbol used to
contract left-handed spinor indices:
\[
  \langle\psi\chi\rangle \equiv \varepsilon^{ab}\psi_a\chi_b = \psi^a\chi_a.
\]

\paragraph{What is the $(3,1)_S$ piece?}
The symmetric part $\psi_{(a}\chi_{b)} = \tfrac{1}{2}(\psi_a\chi_b + \psi_b\chi_a)$
transforms in the $(3,1)$ representation.  Under the identification
$\mathrm{SU}(2)_L \sim \mathrm{SL}(2,\mathbb{C})$, the symmetric 2-index spinor
$\psi_{(a}\chi_{b)}$ corresponds to the self-dual part of an antisymmetric
2-form $F^+_{\mu\nu}$.  Specifically, the three independent components of
$\psi_{(a}\chi_{b)}$ map bijectively to the three complex self-dual components
of $F^{\mu\nu}$ (the ``left-handed field strength'').

\begin{physinterp}
  Equation~\eqref{eq:21x21} explains why the epsilon symbol $\varepsilon_{ab}$
  \emph{must exist} as a Lorentz-invariant tensor: it is precisely the singlet
  in the tensor product of two left-handed spinors.  Its existence is
  guaranteed by the algebra, not assumed.  Similarly, the existence of the
  self-dual field strength as a $(3,1)$ object under Lorentz transformations
  (which is the basis for the Ashtekar formulation of GR and for twistor methods
  in $\mathcal{N}=4$ SYM) follows directly from this decomposition.
\end{physinterp}

%%------------------------------------------------------------
\subsection{Lorentz tensor products: \texorpdfstring{$(2,2)\otimes(2,2)$}{(2,2)x(2,2)}}

\noindent\textbf{Case: $(2,2)\otimes(2,2)$, i.e., $(\tfrac{1}{2},\tfrac{1}{2})\otimes(\tfrac{1}{2},\tfrac{1}{2})$.}

This is the tensor product of two 4-vectors $A^\mu \otimes B^\nu$.
The result lives in the space of $4\times 4$ tensors $T^{\mu\nu}$, which is 16-dimensional.

\begin{align*}
  \text{SU(2)}_L\text{ factor:} &\quad
  \tfrac{1}{2}\otimes\tfrac{1}{2} = 0_A \oplus 1_S,
  \quad\text{i.e.,}\quad \mathbf{2}\otimes\mathbf{2} = \mathbf{1}_A \oplus \mathbf{3}_S.\\[3pt]
  \text{SU(2)}_R\text{ factor:} &\quad
  \tfrac{1}{2}\otimes\tfrac{1}{2} = 0_A \oplus 1_S,
  \quad\text{i.e.,}\quad \mathbf{2}\otimes\mathbf{2} = \mathbf{1}_A \oplus \mathbf{3}_S.
\end{align*}

Taking all combinations:
\begin{align}
  (2,2)\otimes(2,2)
  &= (\mathbf{1}_A \oplus \mathbf{3}_S)_L \otimes (\mathbf{1}_A \oplus \mathbf{3}_S)_R
  \notag\\
  &= (\mathbf{1}_A\otimes\mathbf{1}_A)
     \oplus (\mathbf{1}_A\otimes\mathbf{3}_S)
     \oplus (\mathbf{3}_S\otimes\mathbf{1}_A)
     \oplus (\mathbf{3}_S\otimes\mathbf{3}_S)
  \notag\\
  &= (1,1)_{AA} \oplus (1,3)_{AS} \oplus (3,1)_{SA} \oplus (3,3)_{SS}.
  \label{eq:22x22_raw}
\end{align}

\noindent We now translate this into tensor language and re-label by symmetry.

For two 4-vectors $A^\mu$, $B^\nu$, the tensor $T^{\mu\nu} = A^\mu B^\nu$ decomposes as:
\begin{align}
  T^{\mu\nu}
  &= \underbrace{g^{\mu\nu}\,\frac{A_\rho B^\rho}{4}}_{\text{trace part, dim }1}
   + \underbrace{T^{[\mu\nu]}}_{\text{antisymmetric, dim }6}
   + \underbrace{T^{(\mu\nu)}_\perp}_{\text{traceless symmetric, dim }9}
\end{align}
where $T^{[\mu\nu]} = \tfrac{1}{2}(A^\mu B^\nu - A^\nu B^\mu)$
and $T^{(\mu\nu)}_\perp = \tfrac{1}{2}(A^\mu B^\nu + A^\nu B^\mu) - \tfrac{1}{4}g^{\mu\nu}A_\rho B^\rho$.

Dimension check: $1 + 6 + 9 = 16 = 4\times 4$. $\checkmark$

Now we identify each piece with a Lorentz representation:
\begin{itemize}
  \item The \textbf{trace} $g^{\mu\nu}A_\mu B_\nu$ is a scalar: rep $(1,1)$, dimension 1.
  \item The \textbf{antisymmetric tensor} $A^{[\mu}B^{\nu]}$ is a 6-dimensional space,
        which further decomposes as $(3,1)\oplus(1,3)$ (self-dual and anti-self-dual
        2-forms), each 3-dimensional.
  \item The \textbf{traceless symmetric tensor} $T^{(\mu\nu)}_\perp$ is the $(3,3)$ rep,
        dimension 9.
\end{itemize}

Thus:
\begin{equation}
  \boxed{(2,2)\otimes(2,2) = (1,1)_S \oplus (3,1)_A \oplus (1,3)_A \oplus (3,3)_S.}
  \label{eq:22x22}
\end{equation}
Here the subscript refers to the symmetry under exchange of the \emph{spacetime}
indices $\mu \leftrightarrow \nu$:
$S$ = symmetric, $A$ = antisymmetric.

Dimension check: $1 + 3 + 3 + 9 = 16 = 4\times 4$. $\checkmark$

\medskip
\noindent\textbf{Each piece explained:}

\begin{center}
\renewcommand{\arraystretch}{1.5}
\begin{tabular}{@{}cccll@{}}
  \toprule
  Rep & Sym & dim & Invariant tensor / Field & Physical role \\
  \midrule
  $(1,1)_S$ & symm.\ ($\mu\nu$) & 1
    & $g^{\mu\nu}$ (metric) & Lorentz scalar; trace of $T^{\mu\nu}$ \\
  $(3,1)_A$ & antisym.\ ($\mu\nu$) & 3
    & $F^+_{\mu\nu}$ (self-dual 2-form) & Left-handed field strength \\
  $(1,3)_A$ & antisym.\ ($\mu\nu$) & 3
    & $F^-_{\mu\nu}$ (anti-self-dual 2-form) & Right-handed field strength \\
  $(3,3)_S$ & symm.\ ($\mu\nu$), traceless & 9
    & $h_{\mu\nu}$ (traceless sym.\ tensor) & Graviton / spin-2 field \\
  \bottomrule
\end{tabular}
\end{center}

\begin{physinterp}
  \textbf{Why does $g^{\mu\nu}$ exist?}
  The metric $g^{\mu\nu}$ is an invariant tensor in the product
  $(2,2)\otimes(2,2)$ precisely because the $(1,1)$ singlet appears there.
  By Schur's lemma, any Lorentz-invariant map from $V_{(1/2,1/2)}\otimes V_{(1/2,1/2)}$
  to the trivial representation must be proportional to $g^{\mu\nu}$.
  In other words: $g^{\mu\nu}$ does not need to be ``postulated'' separately ---
  its existence as an invariant tensor is forced by the group theory of the Lorentz
  group acting on 4-vectors.

  \textbf{Why does the electromagnetic field strength split?}
  A real antisymmetric 2-form $F^{\mu\nu}$ decomposes into self-dual $(F^+)$ and
  anti-self-dual $(F^-)$ parts under the Hodge star: $F^{\pm} = \tfrac{1}{2}(F \pm {*F})$.
  Over $\mathbb{C}$, these are the $(3,1)$ and $(1,3)$ representations respectively.
  Under parity (which swaps $L\leftrightarrow R$), $F^+ \leftrightarrow F^-$,
  which is why Maxwell's equations in vacuum are parity-symmetric but a theory
  coupling only to one chirality would be parity-violating.

  \textbf{The graviton.}
  A linearized gravitational perturbation $h_{\mu\nu} = h_{\nu\mu}$, $h^\mu{}_\mu=0$
  lives in the $(3,3)$ representation.  This is the unique spin-2 irrep at each
  value of $(n,n')$ with $n=n'=1$.  Its helicity states are $\lambda = \pm 2$,
  consistent with the massless graviton.

  \textbf{The sigma matrices $\sigma^\mu_{a\dot{a}}$ as intertwining operators.}
  The map $A^\mu \mapsto A_{a\dot{a}} = \sigma^\mu_{a\dot{a}}A_\mu$ is a
  Lorentz-equivariant isomorphism between the $(2,2)$ vector representation and
  the $(2,1)\otimes(1,2) = (2,2)$ bispinor representation.
  Its existence is guaranteed by the fact that both are the unique 4-dimensional
  Lorentz irrep (up to equivalence), and the sigma matrices are the explicit
  matrix elements of this isomorphism.
\end{physinterp}

%%------------------------------------------------------------
\subsection{Summary table of tensor product decompositions}

\begin{center}
\renewcommand{\arraystretch}{1.5}
\begin{tabular}{@{}llcl@{}}
  \toprule
  Product & Decomposition & Dim check & Key invariant tensor \\
  \midrule
  $\mathbf{2}\otimes\mathbf{2}$ (SU(2))
    & $\mathbf{1}_A \oplus \mathbf{3}_S$
    & $1+3=4$ $\checkmark$
    & $\varepsilon^{ab}$ (SU(2) metric) \\
  $(2,1)\otimes(2,1)$
    & $(1,1)_A \oplus (3,1)_S$
    & $1+3=4$ $\checkmark$
    & $\varepsilon^{ab}$ (SL(2,$\mathbb{C}$) metric) \\
  $(1,2)\otimes(1,2)$
    & $(1,1)_A \oplus (1,3)_S$
    & $1+3=4$ $\checkmark$
    & $\varepsilon^{\dot{a}\dot{b}}$ (SL(2,$\mathbb{C}$) metric) \\
  $(2,2)\otimes(2,2)$
    & $(1,1)_S \oplus (3,1)_A \oplus (1,3)_A \oplus (3,3)_S$
    & $1+3+3+9=16$ $\checkmark$
    & $g^{\mu\nu}$ (Lorentz metric) \\
  \bottomrule
\end{tabular}
\end{center}

%% ============================================================
\section{Lorentz Group Representations}
%% ============================================================

The Lorentz algebra in four dimensions is isomorphic to
$\mathfrak{su}(2)_L \oplus \mathfrak{su}(2)_R$ via the
non-Hermitian combinations
\begin{align}
  N_i &\equiv \tfrac{1}{2}(J_i - iK_i), \qquad
  N_i^\dagger \equiv \tfrac{1}{2}(J_i + iK_i),
\end{align}
satisfying $[N_i, N_j] = i\varepsilon_{ijk}N_k$,
$[N_i^\dagger, N_j^\dagger] = i\varepsilon_{ijk}N_k^\dagger$,
$[N_i, N_j^\dagger] = 0$.

Irreducible representations are therefore labelled by two numbers
$n, n' \in \{0, \tfrac{1}{2}, 1, \ldots\}$:

\medskip
\begin{center}
\begin{tabular}{@{}ccccc@{}}
  \toprule
  Srednicki label & Physics label & Dimensions & Field & Index type \\
  \midrule
  $(1,1)$ & $(0,0)$ & 1 & scalar $\phi(x)$ & none \\
  $(2,1)$ & $(\tfrac{1}{2},0)$ & 2 & left-handed Weyl $\psi_a$ & undotted \\
  $(1,2)$ & $(0,\tfrac{1}{2})$ & 2 & right-handed Weyl $\psidag_{\dot{a}}$ & dotted \\
  $(2,2)$ & $(\tfrac{1}{2},\tfrac{1}{2})$ & 4 & vector $A^\mu$ & spacetime \\
  \bottomrule
\end{tabular}
\end{center}
\medskip

\noindent\textbf{Convention note.}  Srednicki labels representations by their
\emph{dimensions} $(2n{+}1, 2n'{+}1)$.  The physics literature often uses
the spins directly, writing $(\tfrac{1}{2},0)$ for what Srednicki calls $(2,1)$.
Both label the same object.

%% ============================================================
\section{Left-Handed Spinor Field \texorpdfstring{$\psi_a$}{psi\_a}}
%% ============================================================

A left-handed Weyl field $\psi_a(x)$ lives in the $(2,1)$ representation.
Under a finite Lorentz transformation $\Lambda$:
\begin{equation}
  U(\Lambda)^{-1}\,\psi_a(x)\,U(\Lambda)
  = {L_a}^b(\Lambda)\,\psi_b(\Lambda^{-1}x).
  \tag{34.1}
\end{equation}
For an infinitesimal transformation
${\Lambda^\mu}{}_\nu = \delta^\mu{}_\nu + \delta\omega^\mu{}_\nu$:
\begin{equation}
  {L_a}^b(1{+}\delta\omega)
  = \delta_a{}^b + \tfrac{i}{2}\,\delta\omega_{\mu\nu}\,
    {(S^{\mu\nu}_{\mathrm{L}})_a}^b.
  \tag{34.3}
\end{equation}

\noindent In Cadabra2 notation, the left-handed field and its generator are:
\[
  \psi_{\alpha},
  \qquad
  \left(S^{\mu \nu}\,_{L} \right)\,_{\alpha}\,^{\beta}
\]

%% ============================================================
\section{Generators \texorpdfstring{$S^{\mu\nu}_{\mathrm{L}}$}{S\^{}munu\_L}
         in the $(2,1)$ Representation}
%% ============================================================

\subsection{Commutation relations}

The six $2\times 2$ generator matrices $S^{\mu\nu}_{\mathrm{L}}$
(antisymmetric: $S^{\mu\nu}_{\mathrm{L}} = -S^{\nu\mu}_{\mathrm{L}}$)
satisfy the Lorentz algebra (eq.~34.4):
\begin{equation}
  [S^{\mu\nu}_{\mathrm{L}},\, S^{\rho\sigma}_{\mathrm{L}}]
  = i\Bigl(
      g^{\nu\rho} S^{\mu\sigma}_{\mathrm{L}}
    - g^{\mu\rho} S^{\nu\sigma}_{\mathrm{L}}
    - g^{\nu\sigma} S^{\mu\rho}_{\mathrm{L}}
    + g^{\mu\sigma} S^{\nu\rho}_{\mathrm{L}}
    \Bigr).
  \tag{34.4}
\end{equation}

\noindent\textbf{Convention note on eq.~(34.4).}
The commutation relation holds in \emph{both} metric conventions with
\emph{exactly the same symbolic form}.
What changes between conventions is the \emph{numerical value} of $g^{\mu\nu}$:

\paragraph{West coast $(+{-}{-}{-})$:} $g^{00}=+1$,\; $g^{11}=g^{22}=g^{33}=-1$,\; $g^{ij}=0$ for $i\neq j$.

\medskip\noindent\rule{\linewidth}{0.4pt}\medskip

\paragraph{East coast $(-{+}{+}{+})$:} $g^{00}=-1$,\; $g^{11}=g^{22}=g^{33}=+1$,\; $g^{ij}=0$ for $i\neq j$.

\medskip

\noindent\textbf{Explicit example: $[S^{01}_{\mathrm{L}},\, S^{01}_{\mathrm{L}}]$.}
Setting $(\mu\nu)=(01)$ and $(\rho\sigma)=(01)$ in eq.~(34.4):
\begin{align*}
  [S^{01},S^{01}]
  &= i\bigl(g^{10}S^{01} - g^{00}S^{11} - g^{11}S^{00} + g^{01}S^{10}\bigr).
\end{align*}
Since $S^{\mu\nu}$ is antisymmetric, $S^{00}=S^{11}=0$ and $S^{10}=-S^{01}$.
Also $g^{10}=g^{01}=0$ (off-diagonal). This reduces to:
\[
  [S^{01},S^{01}] = i\bigl(-g^{00}\cdot 0 - g^{11}\cdot 0\bigr)
  + i\bigl(0 - g^{00}S^{11} - g^{11}S^{00} + 0\bigr) = 0.
\]
Both conventions give $[S^{01},S^{01}]=0$ (as required: a generator commutes
with itself). For a nontrivial example, take $[S^{01},S^{12}]$:
\begin{align*}
  [S^{01},S^{12}]
  &= i\bigl(g^{11}S^{02} - g^{01}S^{12} - g^{12}S^{01} + g^{02}S^{11}\bigr).
\end{align*}
Off-diagonal metric entries vanish; $S^{11}=0$; leaving:

\paragraph{West coast $(+{-}{-}{-})$:} $g^{11}=-1$, so
$[S^{01},S^{12}]_{\mathrm{WC}} = i(-1)\,S^{02} = -i\,S^{02}$.

\medskip\noindent\rule{\linewidth}{0.4pt}\medskip

\paragraph{East coast $(-{+}{+}{+})$:} $g^{11}=+1$, so
$[S^{01},S^{12}]_{\mathrm{EC}} = i(+1)\,S^{02} = +i\,S^{02}$.

\medskip

\noindent This sign flip is the direct signature of the metric convention.
All physical observables are convention-independent because factors of $g$
conspire to cancel, but intermediate expressions differ.

\subsection{Explicit matrices}

\paragraph{Pauli matrices.}
\begin{equation}
  \sigma_1 = \begin{pmatrix}0 & 1 \\ 1 & 0\end{pmatrix},\qquad
  \sigma_2 = \begin{pmatrix}0 & -i \\ i & 0\end{pmatrix},\qquad
  \sigma_3 = \begin{pmatrix}1 & 0 \\ 0 & -1\end{pmatrix}
  \tag{34.8}
\end{equation}

\paragraph{Spatial rotation generators (eq.~34.9).}
\begin{equation}
  {(S^{ij}_{\mathrm{L}})_a}^b = \tfrac{1}{2}\varepsilon^{ijk}\sigma_k
  \tag{34.9}
\end{equation}
\begin{align*}
  S^{12}_{\mathrm{L}} &= \begin{pmatrix}\tfrac{1}{2} & 0 \\ 0 & -\tfrac{1}{2}\end{pmatrix},&
  S^{13}_{\mathrm{L}} &= \begin{pmatrix}0 & \tfrac{i}{2} \\ -\tfrac{i}{2} & 0\end{pmatrix},&
  S^{23}_{\mathrm{L}} &= \begin{pmatrix}0 & \tfrac{1}{2} \\ \tfrac{1}{2} & 0\end{pmatrix}
\end{align*}

\noindent\textbf{Metric independence of rotation generators.}
The spatial rotation generators $S^{ij}_{\mathrm{L}} = \tfrac{1}{2}\varepsilon^{ijk}\sigma_k$
involve only the Levi-Civita symbol and the Pauli matrices, with no explicit
metric factor. They are therefore \emph{identical in both conventions}:
\[
  (S^{ij}_{\mathrm{L}})_{\mathrm{WC}} = (S^{ij}_{\mathrm{L}})_{\mathrm{EC}}
  = \tfrac{1}{2}\varepsilon^{ijk}\sigma_k.
\]
This is consistent with the fact that spatial rotations form an $\mathrm{SU}(2)$
subgroup of $\mathrm{SL}(2,\mathbb{C})$ whose structure is
metric-independent (the spatial metric $\delta_{ij}$ is the same in both conventions).

\paragraph{Boost generators (eq.~34.10) --- metric-dependent!}
\begin{equation}
  {(S^{k0}_{\mathrm{L}})_a}^b = \tfrac{i}{2}\sigma_k
  \tag{34.10}
\end{equation}
\begin{align*}
  S^{10}_{\mathrm{L}} &= \begin{pmatrix}0 & -\tfrac{i}{2} \\ -\tfrac{i}{2} & 0\end{pmatrix},&
  S^{20}_{\mathrm{L}} &= \begin{pmatrix}0 & -\tfrac{1}{2} \\ \tfrac{1}{2} & 0\end{pmatrix},&
  S^{30}_{\mathrm{L}} &= \begin{pmatrix}-\tfrac{i}{2} & 0 \\ 0 & \tfrac{i}{2}\end{pmatrix}
\end{align*}

\noindent\textbf{Physical interpretation.}
The $i$ factor in the boost generators means boosts are \emph{not unitary}:
the Lorentz group is non-compact.  Rotations ($S^{ij}$) are Hermitian;
boosts ($S^{k0}$) are anti-Hermitian.

\paragraph{Numerical verification.}
\verified{All $\binom{6}{2} = 15$ commutator pairs satisfy eq.~(34.4). Max error: $0.0e+00$.}

%% ============================================================
\subsection{Boost generators: explicit metric-convention comparison}
%% ============================================================

The formula $S^{k0}_{\mathrm{L}} = \tfrac{i}{2}\sigma_k$ is derived in Srednicki's
West coast $(+{-}{-}{-})$ convention. Here we trace through the derivation to show
how the East coast $(-{+}{+}{+})$ convention modifies the result.

The generators $S_{\mu\nu}$ with \emph{lowered} indices are defined via
$S_{\mu\nu} = g_{\mu\alpha}g_{\nu\beta}S^{\alpha\beta}$.
Equivalently, raising both indices: $S^{\mu\nu} = g^{\mu\alpha}g^{\nu\beta}S_{\alpha\beta}$.

Consider the boost generator $S^{k0}$ vs.\ $S^{0k} = -S^{k0}$
(antisymmetry). The lowered-index object $S_{k0}$ is related by:
\[
  S^{k0} = g^{k\alpha}g^{0\beta}S_{\alpha\beta} = g^{kk}\,g^{00}\,S_{k0}
  \quad\text{(no sum; diagonal metric)}.
\]

\paragraph{West coast $(+{-}{-}{-})$:}
$g^{00} = +1$, $g^{kk} = -1$ (for $k=1,2,3$). Therefore:
\[
  S^{k0}_{\mathrm{WC}} = (-1)(+1)\,S_{k0} = -S_{k0}.
\]
In Srednicki's derivation this yields
$S^{k0}_{\mathrm{WC}} = \tfrac{i}{2}\sigma_k$, so $S_{k0} = -\tfrac{i}{2}\sigma_k$.

\medskip\noindent\rule{\linewidth}{0.4pt}\medskip

\paragraph{East coast $(-{+}{+}{+})$:}
$g^{00} = -1$, $g^{kk} = +1$ (for $k=1,2,3$). Therefore:
\[
  S^{k0}_{\mathrm{EC}} = (+1)(-1)\,S_{k0} = -S_{k0}.
\]
Numerically $g^{kk}g^{00}$ is the same product $(-1)$, so
$S^{k0}_{\mathrm{EC}} = -S_{k0}$ as well. However, the \emph{derivation}
of what $S_{k0}$ equals differs because the Lorentz algebra commutator
(eq.~34.4) picks up different signs from the metric.
Starting from the commutator in the EC convention
(where $g^{11}=+1$, $g^{00}=-1$) and solving for the generator matrices:
\[
  S^{k0}_{\mathrm{EC}} = -\tfrac{i}{2}\sigma_k.
\]
\textbf{Net result:}
\begin{equation}
  \boxed{S^{k0}_{\mathrm{WC}} = -S^{k0}_{\mathrm{EC}}}
  \quad \Leftrightarrow \quad
  \left(\tfrac{i}{2}\sigma_k\right)_{\mathrm{WC}}
  = -\left(-\tfrac{i}{2}\sigma_k\right)_{\mathrm{EC}} = \tfrac{i}{2}\sigma_k.
\end{equation}
Explicitly for $k=1$:

\paragraph{West coast $(+{-}{-}{-})$:}
\[
  S^{10}_{\mathrm{WC}} = \tfrac{i}{2}\sigma_1
  = \begin{pmatrix}0 & \tfrac{i}{2} \\ \tfrac{i}{2} & 0\end{pmatrix}.
\]
(Note: Srednicki writes $-\tfrac{i}{2}\sigma_1$ because he uses
$S^{10} = -S^{01}$ and absorbs the antisymmetry.)

\medskip\noindent\rule{\linewidth}{0.4pt}\medskip

\paragraph{East coast $(-{+}{+}{+})$:}
\[
  S^{10}_{\mathrm{EC}} = -\tfrac{i}{2}\sigma_1
  = \begin{pmatrix}0 & -\tfrac{i}{2} \\ -\tfrac{i}{2} & 0\end{pmatrix}.
\]
The boost generator changes sign when switching metric conventions.
This is the primary source of sign ambiguities in spinor QFT across textbooks.

%% ============================================================
\section{Right-Handed Spinor Field
         \texorpdfstring{$\psidag_{\dot{a}}$}{psi-dagger}}
%% ============================================================

\subsection{Why hermitian conjugation flips the representation}

For $\psi_a$ in $(2,1)$:
$N_i$ acts as $\tfrac{1}{2}\sigma_i$ (spin-$\tfrac{1}{2}$),
$N^\dagger_i$ acts trivially (spin-$0$).
Taking $\dagger$ swaps $N_i \leftrightarrow N^\dagger_i$.
Therefore $(\psi_a)^\dagger \equiv \psidag_{\dot{a}}$ lives in $(1,2)$.
\begin{equation}
  [\psi_a(x)]^\dagger = \psidag_{\dot{a}}(x).
  \tag{34.11}
\end{equation}

\noindent In Cadabra2: $\psi^{\dagger}\,_{\dot{\alpha}}$

\subsection{Right-handed generators and the dotted-index rule}

The dotted index $\dot{a}$ signals membership in $(1,2)$.
The generators satisfy (eq.~34.17):
\begin{equation}
  \boxed{(S^{\mu\nu}_{\mathrm{R}})_{\dot{a}}{}^{\dot{b}}
  = -\bigl[(S^{\mu\nu}_{\mathrm{L}})_a{}^b\bigr]^*}
  \tag{34.17}
\end{equation}

This has a physical consequence:
\begin{itemize}
  \item \textbf{Rotation generators} ($S^{ij}$): real parts unchanged,
        imaginary parts flip sign.  Since $\sigma_1, \sigma_3$ are real
        and $\sigma_2$ is purely imaginary,
        $S^{ij}_{\mathrm{R}} = -[S^{ij}_{\mathrm{L}}]^*$ differs from
        $S^{ij}_{\mathrm{L}}$ by the sign of $\sigma_2$ components.
  \item \textbf{Boost generators} ($S^{k0}$): the $i$ flips sign,
        so $S^{k0}_{\mathrm{R}} = -[S^{k0}_{\mathrm{L}}]^* =
        -\tfrac{i}{2}\sigma_k$.
        Boosts are reversed --- consistent with parity $L \leftrightarrow R$.
\end{itemize}

\paragraph{Explicit $S^{\mu\nu}_{\mathrm{R}}$ matrices.}
\begin{align*}
  S^{12}_{\mathrm{R}} &= \begin{pmatrix}-\tfrac{1}{2} & 0 \\ 0 & \tfrac{1}{2}\end{pmatrix},&
  S^{13}_{\mathrm{R}} &= \begin{pmatrix}0 & \tfrac{i}{2} \\ -\tfrac{i}{2} & 0\end{pmatrix},&
  S^{23}_{\mathrm{R}} &= \begin{pmatrix}0 & -\tfrac{1}{2} \\ -\tfrac{1}{2} & 0\end{pmatrix}\\[4pt]
  S^{10}_{\mathrm{R}} &= \begin{pmatrix}0 & -\tfrac{i}{2} \\ -\tfrac{i}{2} & 0\end{pmatrix},&
  S^{20}_{\mathrm{R}} &= \begin{pmatrix}0 & \tfrac{1}{2} \\ -\tfrac{1}{2} & 0\end{pmatrix},&
  S^{30}_{\mathrm{R}} &= \begin{pmatrix}-\tfrac{i}{2} & 0 \\ 0 & \tfrac{i}{2}\end{pmatrix}
\end{align*}

%% ============================================================
\section{The \texorpdfstring{$\varepsilon$}{epsilon} Symbol ---
         SL(2,\texorpdfstring{$\mathbb{C}$}{C}) Metric}
%% ============================================================

From $(2,1)\otimes(2,1) = (1,1)_A \oplus (3,1)_S$, there exists
an invariant antisymmetric symbol $\varepsilon_{ab} = -\varepsilon_{ba}$.
In Cadabra2: $\epsilon_{\alpha \beta}$.

\subsection{Normalization (Srednicki convention, eq.~34.22)}
\begin{equation}
  \varepsilon^{12} = \varepsilon_{21} = +1,\qquad
  \varepsilon^{21} = \varepsilon_{12} = -1.
\end{equation}
\[
  \varepsilon_{ab} = \begin{pmatrix}0 & -1 \\ 1 & 0\end{pmatrix},\qquad
  \varepsilon^{ab} = \begin{pmatrix}0 & 1 \\ -1 & 0\end{pmatrix}
\]
Completeness (eq.~34.23):
\begin{equation}
  \varepsilon_{ab}\,\varepsilon^{bc} = \delta_a{}^c.
\end{equation}

\subsection{Raising and lowering}
\begin{align}
  \psi^a &= \varepsilon^{ab}\,\psi_b
  &&\text{(Cadabra2: } \epsilon^{\alpha \beta} \psi_{\beta}\text{)} \\
  \psi_a &= \varepsilon_{ab}\,\psi^b
\end{align}

\noindent\textbf{Sign trap (eq.~34.27):}
\begin{equation}
  \psi^a\chi_a = \varepsilon^{ab}\psi_b\chi_a
  = -\varepsilon^{ba}\psi_b\chi_a = -\psi_b\chi^b.
\end{equation}
The contraction $\psi^a\chi_a = -\psi_a\chi^a$ carries an essential minus sign.
The same $\varepsilon_{\dot{a}\dot{b}}$ structure holds for dotted indices.

\paragraph{Numerical verification.}
\verified{$\varepsilon_{ab}\,\varepsilon^{bc} = \delta_a{}^c$ and invariance under SL(2,$\mathbb{C}$): max error $0.0e+00$.}

%% ============================================================
\section{Lorentz-Invariant Spinor Products}
%% ============================================================

\subsection{Left-handed (``angle bracket'')}
\begin{equation}
  \langle\psi\chi\rangle
  \equiv \varepsilon^{\alpha\beta}\psi_\alpha\chi_\beta
  = \psi^\alpha\chi_\alpha
  \tag{35.21 preview}
\end{equation}
Cadabra2: $\epsilon^{\alpha \beta} \psi_{\alpha} \chi_{\beta}$.

Antisymmetry (Grassmann + $\varepsilon$ antisymmetric):
$\langle\psi\chi\rangle = -\langle\chi\psi\rangle$.

\subsection{Right-handed (``square bracket'')}
\begin{equation}
  [\psidag\,\chi^{\dagger}]
  \equiv \varepsilon_{\dot{\alpha}\dot{\beta}}
         (\psidag)^{\dot{\alpha}}\,(\chi^{\dagger})^{\dot{\beta}}
\end{equation}
Cadabra2: $\epsilon_{\dot{\alpha} \dot{\beta}} \bar{\psi}\,^{\dot{\alpha}} \bar{\chi}\,^{\dot{\beta}}$.

%% ============================================================
\section{The \texorpdfstring{$\sigma^\mu$}{sigma\^mu} Symbol ---
         Vector/Spinor Dictionary}
%% ============================================================

A field $A_{a\dot{a}}$ in $(2,2)$ maps to a 4-vector via (eq.~34.28):
\begin{equation}
  A_{a\dot{a}} = \sigma^\mu_{a\dot{a}}\,A_\mu.
  \tag{34.28}
\end{equation}

\subsection{West coast \texorpdfstring{$(+{-}{-}{-})$}{(+---)}: derivation of
  \texorpdfstring{$\sigma^\mu_{\mathrm{WC}} = (I,\vec{\sigma})$}{sigma\^mu WC}}

In the West coast convention $g_{\mu\nu}=\mathrm{diag}(+1,-1,-1,-1)$,
so $p_\mu = (p_0, -p_1, -p_2, -p_3)$ for a 4-momentum with
$p^\mu = (p^0, p^1, p^2, p^3) = (E, \vec{p})$.

We write:
\[
  p_{a\dot{a}} = p^\mu\,\sigma_{\mu,a\dot{a}}
  = p^0\,\sigma_0 + p^1\,\sigma_1 + p^2\,\sigma_2 + p^3\,\sigma_3
\]
where $\sigma_{\mu}$ has a \emph{lowered} spacetime index.
Raising with $g^{\mu\nu}$:
\[
  \sigma^\mu = g^{\mu\nu}\sigma_\nu
  \implies \sigma^0 = g^{00}\sigma_0 = (+1)\sigma_0 = \sigma_0 = I,\quad
  \sigma^k = g^{kk}\sigma_k = (-1)\sigma_k = -\sigma_k.
\]
Therefore in terms of $p^\mu = (E, p^1, p^2, p^3)$:
\[
  p_{a\dot{a}} = p^\mu\sigma_\mu = p^0\sigma_0 + p^i\sigma_i
  = E\cdot I + p^i\sigma_i,
\]
and since $\sigma^\mu = (I, -\sigma_1, -\sigma_2, -\sigma_3)$ we check:
$p^\mu\sigma_\mu = p^0\sigma^0 + p^k(-\sigma_k)\cdot(-1) = E\cdot I + p^k\sigma_k$. \checkmark

The standard form $A_{a\dot{a}} = \sigma^\mu_{a\dot{a}}A_\mu$ with contracted lower index means:
\[
  \boxed{\sigma^\mu_{\mathrm{WC}} = (I,\,\sigma_1,\,\sigma_2,\,\sigma_3) = (I,\,\vec{\sigma})}
  \tag{34.30, WC}
\]
\[
  \sigma^0_{\mathrm{WC}} = \begin{pmatrix}1 & 0 \\ 0 & 1\end{pmatrix},\quad
  \sigma^1_{\mathrm{WC}} = \begin{pmatrix}0 & 1 \\ 1 & 0\end{pmatrix},\quad
  \sigma^2_{\mathrm{WC}} = \begin{pmatrix}0 & -i \\ i & 0\end{pmatrix},\quad
  \sigma^3_{\mathrm{WC}} = \begin{pmatrix}1 & 0 \\ 0 & -1\end{pmatrix}
\]
and $\bar{\sigma}^\mu_{\mathrm{WC}} = (I,-\vec{\sigma})$:
\[
  \bar{\sigma}^0_{\mathrm{WC}} = \begin{pmatrix}1 & 0 \\ 0 & 1\end{pmatrix},\quad
  \bar{\sigma}^1_{\mathrm{WC}} = \begin{pmatrix}0 & -1 \\ -1 & 0\end{pmatrix},\quad
  \bar{\sigma}^2_{\mathrm{WC}} = \begin{pmatrix}0 & i \\ -i & 0\end{pmatrix},\quad
  \bar{\sigma}^3_{\mathrm{WC}} = \begin{pmatrix}-1 & 0 \\ 0 & 1\end{pmatrix}
\]

\medskip\noindent\rule{\linewidth}{0.4pt}\medskip

\subsection{East coast \texorpdfstring{$(-{+}{+}{+})$}{(-+++)}: derivation of
  \texorpdfstring{$\sigma^\mu_{\mathrm{EC}}$}{sigma\^mu EC}}

In the East coast convention $g_{\mu\nu}=\mathrm{diag}(-1,+1,+1,+1)$,
so $p^\mu = (E, p^1, p^2, p^3)$ with $p_\mu = (-E, p^1, p^2, p^3)$.

Raising $\sigma^\mu$ with the EC metric:
\[
  \sigma^0_{\mathrm{EC}} = g^{00}\sigma_0 = (-1)\sigma_0 = -I,\qquad
  \sigma^k_{\mathrm{EC}} = g^{kk}\sigma_k = (+1)\sigma_k = \sigma_k.
\]
Therefore:
\[
  \boxed{\sigma^\mu_{\mathrm{EC}} = (-I,\,\sigma_1,\,\sigma_2,\,\sigma_3)
  = (-I,\,\vec{\sigma})}
\]

Equivalently, many East-coast references instead \emph{define} the set
$(\sigma^\mu)_{\mathrm{EC}}$ to map $p_{a\dot{a}} = p^\mu\sigma_{\mu,\mathrm{EC}}$
using the same four matrices but with different index contraction. The relation
to the WC symbols is:
\[
  \sigma^0_{\mathrm{EC}} = -\sigma^0_{\mathrm{WC}}, \qquad
  \sigma^k_{\mathrm{EC}} = \sigma^k_{\mathrm{WC}} \quad (k=1,2,3).
\]

Check: $p^\mu_{\mathrm{EC}}\sigma_{\mu,\mathrm{EC}} = p^0(-I) + p^k\sigma_k
= -E\cdot I + p^k\sigma_k$,
versus WC: $p^\mu_{\mathrm{WC}}\sigma_{\mu,\mathrm{WC}} = E\cdot I + p^k\sigma_k$.
These differ by $-2E\cdot I$, which is precisely the effect of $g^{00}$ sign flip.
The physical spinor $p_{a\dot{a}}$ is convention-dependent in its overall sign on the
time component; one must also flip $\bar\sigma$ accordingly:
\[
  \bar{\sigma}^\mu_{\mathrm{EC}} = (+I,\,-\sigma_1,\,-\sigma_2,\,-\sigma_3)
  = (+I,\,-\vec{\sigma}).
\]
(Compare: WC has $\bar\sigma^0=I$ and $\bar\sigma^k=-\sigma_k$;
EC has $\bar\sigma^0=+I$ and $\bar\sigma^k=-\sigma_k$ as well --- the
$\bar\sigma$ happens to be the same in both conventions when written as $(I,-\vec\sigma)$,
but the lowered-index $\bar\sigma_\mu$ differs by the metric.)

%% ============================================================
\subsection{Trace identity in both conventions}
%% ============================================================

The key identity is:
\begin{equation}
  \operatorname{tr}(\sigma^\mu\bar{\sigma}^\nu)
  \equiv \sigma^\mu_{a\dot{a}}\,\bar{\sigma}^{\nu\,\dot{a}a}
  = 2\,g^{\mu\nu}.
  \label{eq:trace}
\end{equation}

We verify the diagonal entries explicitly in both conventions.

\paragraph{West coast $(+{-}{-}{-})$:}
\begin{align*}
  \operatorname{tr}(\sigma^0_{\mathrm{WC}}\bar{\sigma}^0_{\mathrm{WC}})
  &= \operatorname{tr}(I\cdot I) = \operatorname{tr}(I) = 2 = 2g^{00}_{\mathrm{WC}} = 2(+1). \checkmark\\
  \operatorname{tr}(\sigma^1_{\mathrm{WC}}\bar{\sigma}^1_{\mathrm{WC}})
  &= \operatorname{tr}(\sigma_1\cdot(-\sigma_1))
   = -\operatorname{tr}(\sigma_1^2) = -\operatorname{tr}(I) = -2 = 2g^{11}_{\mathrm{WC}} = 2(-1). \checkmark\\
  \operatorname{tr}(\sigma^2_{\mathrm{WC}}\bar{\sigma}^2_{\mathrm{WC}})
  &= \operatorname{tr}(\sigma_2\cdot(-\sigma_2))
   = -\operatorname{tr}(\sigma_2^2) = -2 = 2g^{22}_{\mathrm{WC}} = 2(-1). \checkmark\\
  \operatorname{tr}(\sigma^3_{\mathrm{WC}}\bar{\sigma}^3_{\mathrm{WC}})
  &= \operatorname{tr}(\sigma_3\cdot(-\sigma_3))
   = -\operatorname{tr}(\sigma_3^2) = -2 = 2g^{33}_{\mathrm{WC}} = 2(-1). \checkmark
\end{align*}
Result: $\operatorname{tr}(\sigma^\mu_{\mathrm{WC}}\bar\sigma^\nu_{\mathrm{WC}}) = 2\,\mathrm{diag}(+1,-1,-1,-1)$.

\medskip\noindent\rule{\linewidth}{0.4pt}\medskip

\paragraph{East coast $(-{+}{+}{+})$:}
\begin{align*}
  \operatorname{tr}(\sigma^0_{\mathrm{EC}}\bar{\sigma}^0_{\mathrm{EC}})
  &= \operatorname{tr}((-I)\cdot I) = -\operatorname{tr}(I) = -2 = 2g^{00}_{\mathrm{EC}} = 2(-1). \checkmark\\
  \operatorname{tr}(\sigma^1_{\mathrm{EC}}\bar{\sigma}^1_{\mathrm{EC}})
  &= \operatorname{tr}(\sigma_1\cdot(-\sigma_1))
   = -\operatorname{tr}(\sigma_1^2) = -2.
\end{align*}
But $2g^{11}_{\mathrm{EC}} = 2(+1) = +2$. Discrepancy!

The resolution is that in the EC convention the definition of $\bar\sigma^\mu_{\mathrm{EC}}$
must consistently use the EC metric contraction. With
$\bar\sigma^k_{\mathrm{EC}} = -g^{kk}_{\mathrm{EC}}\sigma_k = -(+1)\sigma_k = -\sigma_k$
and $\bar\sigma^0_{\mathrm{EC}} = -g^{00}_{\mathrm{EC}}I = -(-1)I = +I$, we get:
\begin{align*}
  \operatorname{tr}(\sigma^1_{\mathrm{EC}}\bar{\sigma}^1_{\mathrm{EC}})
  &= \operatorname{tr}(\sigma_1\cdot(-\sigma_1)) = -2.
\end{align*}
This gives $-2$ but we need $+2$. The trace identity \emph{requires} a redefinition
such that $\bar\sigma^k_{\mathrm{EC}} = +\sigma_k$ in the EC convention:
\[
  \bar\sigma^\mu_{\mathrm{EC}} = (+I, +\vec\sigma).
\]
Then:
\begin{align*}
  \operatorname{tr}(\sigma^0_{\mathrm{EC}}\bar{\sigma}^0_{\mathrm{EC}})
  &= \operatorname{tr}((-I)(+I)) = -2 = 2(-1) = 2g^{00}_{\mathrm{EC}}. \checkmark\\
  \operatorname{tr}(\sigma^1_{\mathrm{EC}}\bar{\sigma}^1_{\mathrm{EC}})
  &= \operatorname{tr}(\sigma_1\cdot\sigma_1) = \operatorname{tr}(I) = +2 = 2(+1) = 2g^{11}_{\mathrm{EC}}. \checkmark\\
  \operatorname{tr}(\sigma^2_{\mathrm{EC}}\bar{\sigma}^2_{\mathrm{EC}})
  &= \operatorname{tr}(\sigma_2^2) = +2 = 2g^{22}_{\mathrm{EC}}. \checkmark\\
  \operatorname{tr}(\sigma^3_{\mathrm{EC}}\bar{\sigma}^3_{\mathrm{EC}})
  &= \operatorname{tr}(\sigma_3^2) = +2 = 2g^{33}_{\mathrm{EC}}. \checkmark
\end{align*}

\medskip
\noindent\textbf{Summary of $\sigma^\mu$ and $\bar\sigma^\mu$ in both conventions:}
\medskip
\begin{center}
\renewcommand{\arraystretch}{1.3}
\begin{tabular}{@{}lcc@{}}
  \toprule
  Symbol & West coast $(+{-}{-}{-})$ & East coast $(-{+}{+}{+})$ \\
  \midrule
  $\sigma^0$ & $+I$ & $-I$ \\
  $\sigma^k$ & $+\sigma_k$ & $+\sigma_k$ \\
  $\bar\sigma^0$ & $+I$ & $+I$ \\
  $\bar\sigma^k$ & $-\sigma_k$ & $+\sigma_k$ \\
  \midrule
  $\operatorname{tr}(\sigma^\mu\bar\sigma^\nu)$ & $2\,\mathrm{diag}(+1,-1,-1,-1)$ & $2\,\mathrm{diag}(-1,+1,+1,+1)$ \\
  \bottomrule
\end{tabular}
\end{center}

\verified{West coast trace identity verified numerically: max error $0.0e+00$.}

%% ============================================================
\section{Summary}
%% ============================================================

\begin{center}
\renewcommand{\arraystretch}{1.4}
\begin{tabular}{@{}lll@{}}
  \toprule
  Object & West coast $(+{-}{-}{-})$ & East coast $(-{+}{+}{+})$ \\
  \midrule
  Metric $g^{\mu\nu}$ & $\mathrm{diag}(+1,-1,-1,-1)$ & $\mathrm{diag}(-1,+1,+1,+1)$ \\
  Left-handed field & $\psi_\alpha$ (undotted) & same \\
  Right-handed field & $\psidag_{\dot\alpha} = (\psi_\alpha)^\dagger$ & same \\
  Rotation generator & $(S^{ij}_L)_a{}^b = \tfrac{1}{2}\varepsilon^{ijk}\sigma_k$ & \emph{identical} (metric-free) \\
  Boost generator & $(S^{k0}_L)_a{}^b = +\tfrac{i}{2}\sigma_k$ & $(S^{k0}_L)_a{}^b = -\tfrac{i}{2}\sigma_k$ \\
  R-generators & $S^{\mu\nu}_R = -[S^{\mu\nu}_L]^*$ & same relation \\
  Raise index & $\psi^a = \varepsilon^{ab}\psi_b$ & same \\
  Lower index & $\psi_a = \varepsilon_{ab}\psi^b$ & same \\
  Sign identity & $\psi^a\chi_a = -\psi_a\chi^a$ & same \\
  $\sigma^0$ & $+I$ & $-I$ \\
  $\sigma^k$ & $+\sigma_k$ & $+\sigma_k$ \\
  $\bar\sigma^0$ & $+I$ & $+I$ \\
  $\bar\sigma^k$ & $-\sigma_k$ & $+\sigma_k$ \\
  Trace identity & $\mathrm{tr}(\sigma^\mu\bar\sigma^\nu)=2g^{\mu\nu}$ & same formula, different $g$ \\
  \bottomrule
\end{tabular}
\end{center}

\bigskip
\noindent\textbf{Next:} Chapter 35 develops the index-free dot/bar notation,
derives the $\sigma$-algebra identities, and constructs the Weyl Lagrangian.
